\begin{scanpage}{0031}{14}
where
\[g=\frac{\sin\frac\pi n}{n}\left(\frac1{\cos\frac\pi n-\cos\frac{\phi-\phi_0}{n}}+\frac1{\cos\frac\pi n-\cos\frac{\phi+\phi_0}{n}}\right).\tag{2.08}\]
In the vicinity of the shadow boundary $(\phi\approx\pi+\phi_0)$ and near the directions of the mirror-reflected rays $(\phi\approx\pi-\phi_0,\ \phi\approx2\alpha-\pi-\phi_0)$ Expressions (2.04)--(2.08) are not valid, since the poles
\[s_1=\sqrt2e^{i\frac\pi4}\sin\frac{\pi-\psi}{2},\quad s_2=\sqrt2e^{i\frac\pi4}\sin\left(n\pi-\frac{\pi+\psi}{2}\right)\]
of the integrand in (2.02) are close to $s=0$ and, consequently, its expansion in a Taylor series loses meaning. Physically, this result means that in the indicated region the diffraction wave does not reduce to plane and cylindrical waves, but has a more complicated character. An asymptotic representation of the function $v(r,\psi)$ in this region was obtained in 1938 by Pauli [22]; here we will present the derivation of the first term of the asymptotic series obtained in [22].

Let us multiply and divide the integrand expression in Equation (2.02) by the quantity
\[\cos\psi+\cos\zeta=i(s^2-is_0^2)\quad\left(s_0^2=2\cos^2\frac\psi2\right)\tag{2.09}\]
and let us expand into a Taylor series in powers of $s$ the function
\[\frac{\cos\psi+\cos\zeta}{\left(\cos\frac\pi n-\cos\frac{\psi+\zeta}{n}\right)\cos\frac\zeta2},\]
which no longer has a pole at the saddle point $(s=0)$ when $\psi=\phi\pm\phi_0=\pi$. Limiting ourselves in this series to the first term, we obtain
\[v(r,\psi)=\frac{\sqrt2}{\pi}\frac{\sin\frac\pi n}{n}\frac{1+\cos\psi}{\cos\frac\pi n-\cos\frac\psi n}e^{i(kr-\frac\pi4)}\int_0^\infty\frac{e^{-krs^2}}{s^2-is_0^2}\,ds.\tag{2.10}\]
\end{scanpage}
